\section{Benchmark agenda: details}


\textbf{Toward comparable numbers.} For a suite to accumulate results across policies and labs, each type needs a \emph{canonical} scalar and predicate rather than a study-specific one: a keep-out radius per hazard class for T1, a link-to-body margin for T2, a presentation angle $\theta$ for T3, a tilt or spill criterion for T4, the ISO/TS 15066 speed-and-separation and power-and-force bounds for T5, and a contact or time-to-collision criterion for T6. Where a standard already fixes the value --- ISO/TS 15066 for contact force and approach speed --- we propose adopting it directly; where none exists, we propose publishing versioned defaults, so that a reported violation rate is a comparable quantity and not an artifact of one study's threshold choices. The fixability ablations should likewise be standardized: the language ablation and the perception ablation are cheap, decisive, and belong in every type's protocol, because they are what convert a raw defect rate into a claim about the \emph{kind} of fix required.

\textbf{Open design decisions} we put to the community: (i) whether T3 (orientation) should be elevated to a full second pillar via a dedicated handover benchmark, since it is the most intuitive instance of the thesis (``hand the knife handle-first''); (ii) whether T5b (contact force) and T2 (swept volume) merge into one ``robot-body physical safety'' sub-type or remain distinct; (iii) which tasks and scenes each dimension should add first --- the benchmark's next round.

